\chapter{1991年全国卷（理）}
\section{单选题}

\begin{problem}
已知 \(\sin \alpha = \frac{4}{5}\)，并且 \(\alpha\) 是第二象限的角，那么 \(\tan \alpha\) 的值等于（\(\quad\)）

\choices
    {\(-\frac{4}{3}\)}
    {\(-\frac{3}{4}\)}
    {\(\frac{3}{4}\)}
    {\(\frac{4}{3}\)}

\answer{A}

\begin{solution}
由 \(\sin^2\alpha+\cos^2\alpha=1\)，得 \(\cos^2\alpha=1-\frac{16}{25}=\frac{9}{25}\). 因为 \(\alpha\) 是第二象限的角，所以 \(\cos\alpha=-\frac{3}{5}\)，从而 \(\tan\alpha=\frac{\sin\alpha}{\cos\alpha}=-\frac{4}{3}\). 故选 A.
\end{solution}
\end{problem}

\begin{problem}
焦点在 \(\left(-1,0\right)\)，顶点在 \(\left(1,0\right)\) 的抛物线方程是（\(\quad\)）

\choices
    {\(y^{2}=8\left(x+1\right)\)}
    {\(y^{2}=-8\left(x+1\right)\)}
    {\(y^{2}=8\left(x-1\right)\)}
    {\(y^{2}=-8\left(x-1\right)\)}

\answer{D}

\begin{solution}
抛物线的焦点和顶点都在 \(x\) 轴上，且焦点 \((-1,0)\) 在顶点 \((1,0)\) 的左侧，所以抛物线开口向左. 顶点到焦点的距离为 \(2\)，故其标准方程为 \(y^2=-4\times2\left(x-1\right)=-8\left(x-1\right)\). 故选 D.
\end{solution}
\end{problem}

\begin{problem}
函数 \(y=\cos^4 x-\sin^4 x\) 的最小正周期是（\(\quad\)）

\choices
    {\(\frac{\pi}{2}\)}
    {\(\pi\)}
    {\(2\pi\)}
    {\(4\pi\)}

\answer{B}

\begin{solution}
利用平方差公式，\(\cos^4x-\sin^4x=\left(\cos^2x-\sin^2x\right)\left(\cos^2x+\sin^2x\right)=\cos2x\). 函数 \(\cos2x\) 的最小正周期为 \(\pi\)，故选 B.
\end{solution}
\end{problem}

\begin{problem}
如果把两条异面直线看成“一对”，那么六棱锥的棱所在的 \(12\) 条直线中，异面直线共有（\(\quad\)）

\choices
    {\(12\) 对}
    {\(24\) 对}
    {\(36\) 对}
    {\(48\) 对}

\answer{B}

\begin{solution}
六棱锥有 \(6\) 条底棱和 \(6\) 条侧棱. 底棱所在直线都在同一个底面内，任取两条侧棱所在直线都经过顶点，因此这两类直线内部均不构成异面直线. 一条底棱所在直线与一条侧棱所在直线共有 \(6\times6=36\) 对，其中侧棱与底棱有公共端点的情形为 \(6\times2=12\) 对，均相交. 对于其余的底棱直线 \(AB\) 和侧棱直线 \(SC\)，若它们共面，则该平面同时包含底面内不共线的 \(A\)、\(B\)、\(C\)，只能是底面，从而包含不在底面内的顶点 \(S\)，矛盾；所以它们既不相交也不共面. 因此异面直线共有 \(36-12=24\) 对，选 B.
\end{solution}
\end{problem}

\begin{problem}
函数 \(y=\sin\left(2x+\frac{5\pi}{2}\right)\) 的图象的一条对称轴的方程是（\(\quad\)）

\choices
    {\(x = -\frac{\pi}{2}\)}
    {\(x = -\frac{\pi}{4}\)}
    {\(x = \frac{\pi}{8}\)}
    {\(x = \frac{5\pi}{4}\)}

\answer{A}

\begin{solution}
因为 \(\sin\left(2x+\frac{5\pi}{2}\right)=\sin\left(2x+\frac{\pi}{2}\right)=\cos2x\)，所以函数图象是 \(y=\cos2x\). 其对称轴为 \(x=\frac{k\pi}{2}\)，其中 \(k\in\Z\)，取 \(k=-1\) 得 \(x=-\frac{\pi}{2}\). 故选 A.
\end{solution}
\end{problem}

\begin{problem}
如果三棱锥 \(S-ABC\) 的底面是不等边三角形，侧面与底面所成的二面角都相等，且顶点 \(S\) 在底面的射影 \(O\) 在 \(\triangle ABC\) 内，那么 \(O\) 是 \(\triangle ABC\) 的（\(\quad\)）

\choices
    {垂心}
    {重心}
    {外心}
    {内心}

\answer{D}

\begin{solution}
设 \(H\) 是点 \(O\) 到底面边 \(BC\) 的垂足. 过 \(SO\) 作垂直于 \(BC\) 的截面，则该截面与底面、侧面 \(SBC\) 的交线分别为 \(OH\)、\(SH\)，所以对应的二面角满足 \(\tan\theta=\frac{SO}{OH}\). 三个侧面与底面所成的二面角相等，且 \(SO>0\)，故点 \(O\) 到三条边的距离相等. 又因 \(O\) 在三角形内部，\(O\) 是 \(\triangle ABC\) 的内心，选 D.
\end{solution}
\end{problem}

\begin{problem}
已知 \(\{a_n\}\) 是等比数列，且 \(a_n>0\)，\(a_2a_4+2a_3a_5+a_4a_6=25\)，那么 \(a_3+a_5\) 的值等于（\(\quad\)）

\choices
    {\(5\)}
    {\(10\)}
    {\(15\)}
    {\(20\)}

\answer{A}

\begin{solution}
设等比数列的公比为 \(q\)，并记 \(a_3=u\). 则 \(a_2a_4=u^2\)，\(a_3a_5=u^2q^2\)，\(a_4a_6=u^2q^4\). 由题意，\(u^2\left(1+2q^2+q^4\right)=u^2\left(1+q^2\right)^2=25\). 因为各项为正数，\(u>0\)，所以 \(u(1+q^2)=5\). 因而 \(a_3+a_5=u+uq^2=5\)，选 A.
\end{solution}
\end{problem}

\begin{problem}
如果圆锥曲线的极坐标方程为 \(\rho=\frac{16}{5-3\cos\theta}\)，那么它的焦点的极坐标为（\(\quad\)）

\choices
    {\(\left(0,0\right)\)，\(\left(6,\pi\right)\)}
    {\(\left(-3,0\right)\)，\(\left(3,0\right)\)}
    {\(\left(0,0\right)\)，\(\left(3,0\right)\)}
    {\(\left(0,0\right)\)，\(\left(6,0\right)\)}

\answer{D}

\begin{solution}
令 \(\rho=\sqrt{x^2+y^2}\)，则 \(\rho\cos\theta=x\). 原极坐标方程化为 \(5\rho=16+3x\). 由于 \(5-3\cos\theta>0\)，两边平方不引入额外点，得 \(25(x^2+y^2)=(16+3x)^2\)，即 \(16(x-3)^2+25y^2=400\)，或 \(\frac{(x-3)^2}{25}+\frac{y^2}{16}=1\). 该椭圆的中心为 \((3,0)\)，半长轴为 \(5\)，半短轴为 \(4\)，焦距的一半为 \(\sqrt{5^2-4^2}=3\). 所以两个焦点为 \((0,0)\)、\((6,0)\)，选 D.
\end{solution}
\end{problem}

\begin{problem}
从 \(4\) 台甲型和 \(5\) 台乙型电视机中任意取出 \(3\) 台，其中至少要有甲型与乙型电视机各 \(1\) 台，则不同的取法共有（\(\quad\)）

\choices
    {\(140\) 种}
    {\(84\) 种}
    {\(70\) 种}
    {\(35\) 种}

\answer{C}

\begin{solution}
从 \(9\) 台电视机中取出 \(3\) 台的总取法数为 \(\mathrm C_9^3\). 其中全取甲型有 \(\mathrm C_4^3\) 种，全取乙型有 \(\mathrm C_5^3\) 种. 因而至少各有一台的取法数为 \(\mathrm C_9^3-\mathrm C_4^3-\mathrm C_5^3=84-4-10=70\)，选 C.
\end{solution}
\end{problem}

\begin{problem}
如果 \(AC<0\) 且 \(BC<0\)，那么直线 \(Ax+By+C=0\) 不通过（\(\quad\)）

\choices
    {第一象限}
    {第二象限}
    {第三象限}
    {第四象限}

\answer{C}

\begin{solution}
由 \(AC<0\)、\(BC<0\) 可知 \(A\)、\(B\) 同号且与 \(C\) 异号. 将方程整体乘以 \(-1\) 后不改变直线，可设 \(A>0\)、\(B>0\)、\(C<0\). 当点在第三象限时，\(x<0\)、\(y<0\)，于是 \(Ax+By+C<0\)，不可能等于 \(0\). 因此直线不通过第三象限，选 C.
\end{solution}
\end{problem}

\begin{problem}
设甲，乙，丙是三个命题. 如果甲是乙的必要条件；丙是乙的充分条件但不是乙的必要条件，那么（\(\quad\)）

\choices
    {丙是甲的充分条件，但不是甲的必要条件}
    {丙是甲的必要条件，但不是甲的充分条件}
    {丙是甲的充要条件}
    {丙不是甲的充分条件，也不是甲的必要条件}

\answer{A}

\begin{solution}
“甲是乙的必要条件”表示 \(\text{乙}\Rightarrow\text{甲}\)；“丙是乙的充分条件”表示 \(\text{丙}\Rightarrow\text{乙}\). 因而 \(\text{丙}\Rightarrow\text{乙}\Rightarrow\text{甲}\)，丙是甲的充分条件. 若甲也是丙的充分条件，则有 \(\text{甲}\Rightarrow\text{丙}\Rightarrow\text{乙}\)，这将推出 \(\text{乙}\Rightarrow\text{丙}\)，与丙不是乙的必要条件矛盾. 所以丙不是甲的必要条件，选 A.
\end{solution}
\end{problem}

\begin{problem}
\(\displaystyle\lim_{n\to\infty}\left[n\left(1-\frac{1}{3}\right)\left(1-\frac{1}{4}\right)\left(1-\frac{1}{5}\right)\cdots\left(1-\frac{1}{n+2}\right)\right]\) 的值等于（\(\quad\)）

\choices
    {\(0\)}
    {\(1\)}
    {\(2\)}
    {\(3\)}

\answer{C}

\begin{solution}
将乘积展开可得
\(n\left(1-\frac{1}{3}\right)\left(1-\frac{1}{4}\right)\cdots\left(1-\frac{1}{n+2}\right)=n\cdot\frac{2}{3}\cdot\frac{3}{4}\cdots\frac{n+1}{n+2}=\frac{2n}{n+2}\). 因此极限为 \(\displaystyle\lim_{n\to\infty}\frac{2n}{n+2}=2\)，选 C.
\end{solution}
\end{problem}

\begin{problem}
如果奇函数 \(f(x)\) 在区间 \([3,7]\) 上是增函数且最小值为 \(5\)，那么 \(f(x)\) 在区间 \([-7,-3]\) 上是（\(\quad\)）

\choices
    {增函数且最小值为 \(-5\)}
    {增函数且最大值为 \(-5\)}
    {减函数且最小值为 \(-5\)}
    {减函数且最大值为 \(-5\)}

\answer{B}

\begin{solution}
任取 \(x_1\)，\(x_2\in[-7,-3]\)，令 \(t_1=-x_1\)、\(t_2=-x_2\)，则 \(t_1\)，\(t_2\in[3,7]\). 当 \(x_1<x_2\) 时，\(t_1>t_2\). 因为 \(f\) 在 \([3,7]\) 上递增，所以 \(f(t_1)>f(t_2)\)，再由奇函数性质得 \(f(x_1)=-f(t_1)<-f(t_2)=f(x_2)\). 因此 \(f\) 在 \([-7,-3]\) 上递增. 又因为 \(f\) 在 \([3,7]\) 上递增，所以其最小值在 \(x=3\) 处取得，即 \(f(3)=5\)，从而 \(f(-3)=-f(3)=-5\). 因此 \(f\) 在 \([-7,-3]\) 上的最大值为 \(-5\)，选 B.
\end{solution}
\end{problem}

\begin{problem}
圆 \(x^{2} + 2x + y^{2} + 4y - 3 = 0\) 上到直线 \(x + y + 1 = 0\) 的距离为 \(\sqrt{2}\) 的点共有（\(\quad\)）

\choices
    {\(1\) 个}
    {\(2\) 个}
    {\(3\) 个}
    {\(4\) 个}

\answer{C}

\begin{solution}
圆的方程化为 \(\left(x+1\right)^2+\left(y+2\right)^2=8\)，圆心为 \((-1,-2)\)，半径为 \(2\sqrt2\). 设所求点到直线的距离为 \(\sqrt2\)，则 \(\lvert x+y+1\rvert=2\)，所以这些点在两条直线 \(x+y=1\) 或 \(x+y=-3\) 上. 直线 \(x+y=1\) 到圆心的距离为 \(2\sqrt{2}\)，与圆相切，有 \(1\) 个交点；直线 \(x+y=-3\) 经过圆心，与圆有 \(2\) 个交点. 共有 \(3\) 个点，选 C.
\end{solution}
\end{problem}

\begin{problem}
设全集为 \(\R\)，\(f(x)=\sin x\)，\(g(x)=\cos x\)，\(M=\{x\mid f(x)\neq 0\}\)，\(N=\{x\mid g(x)\neq 0\}\)，那么集合 \(\{x\mid f(x)g(x)=0\}\) 等于（\(\quad\)）

\choices
    {\(\overline{M} \cap \overline{N}\)}
    {\(\overline{M} \cup N\)}
    {\(M \cup \overline{N}\)}
    {\(\overline{M} \cup \overline{N}\)}

\answer{D}

\begin{solution}
由定义，\(\overline M=\{x\mid\sin x=0\}\)，\(\overline N=\{x\mid\cos x=0\}\). 条件 \(f(x)g(x)=0\) 等价于 \(\sin x=0\) 或 \(\cos x=0\)，所以 \(\{x\mid f(x)g(x)=0\}=\overline M\cup\overline N\)，选 D.
\end{solution}
\end{problem}

\section{填空题}

\begin{problem}
\(\arctan \frac{1}{3}+\arctan \frac{1}{2}\) 的值是\fillinblank{}.

\answer{\(\frac{\pi}{4}\)}

\begin{solution}
设 \(u=\arctan\frac{1}{3}+\arctan\frac{1}{2}\)，则 \(0<u<\pi\). 由正切和角公式，\(\tan u=\frac{\frac13+\frac12}{1-\frac16}=1>0\)，且分母为正说明 \(\cos u>0\)，所以 \(u\in\left(0,\frac{\pi}{2}\right)\). 因而 \(u=\frac{\pi}{4}\).
\end{solution}
\end{problem}

\begin{problem}
不等式 \(6^{x^2+x-2}<1\) 的解集是\fillinblank{}.

\answer{\(\left(-2,1\right)\)}

\begin{solution}
因为底数 \(6>1\)，所以 \(6^{x^2+x-2}<1\) 等价于 \(x^2+x-2<0\). 因式分解得 \((x+2)(x-1)<0\)，故解集为 \(\left(-2,1\right)\).
\end{solution}
\end{problem}

\begin{problem}
已知正三棱台上底面边长为 \(2\)，下底面边长为 \(4\)，且侧棱与底面所成的角是 \(45^\circ\)，那么这个正三棱台的体积等于\fillinblank{}.

\answer{\(\frac{14}{3}\)}

\begin{solution}
设该正三棱台的高为 \(h\). 正三棱台上下底面的中心在同一条高线上，对应顶点连线在底面上的投影长等于两底面外接圆半径之差. 正三角形边长分别为 \(2\)、\(4\) 时，外接圆半径分别为 \(\frac{2}{\sqrt3}\)、\(\frac{4}{\sqrt3}\)，故投影长为 \(\frac{2}{\sqrt3}\). 在由侧棱、其投影和高组成的直角三角形中，\(\tan45^\circ=\frac{h}{2/\sqrt3}=1\)，所以 \(h=\frac{2}{\sqrt3}\). 上、下底面积分别为 \(\sqrt3\)、\(4\sqrt3\)，故正棱台体积为
\[
V=\frac{h}{3}\left(S_1+S_2+\sqrt{S_1S_2}\right)=\frac{2/\sqrt3}{3}\left(\sqrt3+4\sqrt3+2\sqrt3\right)=\frac{14}{3}
\]
\end{solution}
\end{problem}

\begin{problem}
\(\left(ax+1\right)^7\) 的展开式中，\(x^3\) 的系数是 \(x^2\) 的系数与 \(x^4\) 的系数的等差中项，若实数 \(a>1\)，那么 \(a=\)\fillinblank{}.

\answer{\(1+\frac{\sqrt{10}}{5}\)}

\begin{solution}
在 \((ax+1)^7\) 的展开式中，\(x^k\) 的系数为 \(\mathrm C_7^k a^k\). 因而 \(x^2\)，\(x^3\)，\(x^4\) 的系数分别为 \(21a^2\)、\(35a^3\)、\(35a^4\). 由等差中项关系，\(2\times35a^3=21a^2+35a^4\). 因为 \(a>1\)，可除以 \(7a^2\)，得到 \(5a^2-10a+3=0\)，所以 \(a=1\pm\frac{\sqrt{10}}{5}\). 结合 \(a>1\)，得 \(a=1+\frac{\sqrt{10}}{5}\).
\end{solution}
\end{problem}

\begin{problem}
在球面上有四个点 \(P\)，\(A\)，\(B\)，\(C\)，如果 \(PA\)，\(PB\)，\(PC\) 两两互相垂直，且 \(PA=PB=PC=a\)，那么这个球面的面积是\fillinblank{}.

\answer{\(3\pi a^2\)}

\begin{solution}
以 \(P\) 为原点，分别以 \(PA\)、\(PB\)、\(PC\) 所在方向为三个坐标轴正方向，则可设 \(P=(0,0,0)\)、\(A=(a,0,0)\)、\(B=(0,a,0)\)、\(C=(0,0,a)\). 设球心为 \(O=(u,v,w)\). 由 \(OP=OA\) 得 \(u^2+v^2+w^2=(u-a)^2+v^2+w^2\)，所以 \(u=\frac a2\). 由 \(OP=OB\) 得 \(u^2+v^2+w^2=u^2+(v-a)^2+w^2\)，所以 \(v=\frac a2\). 由 \(OP=OC\) 得 \(u^2+v^2+w^2=u^2+v^2+(w-a)^2\)，所以 \(w=\frac a2\). 因而球半径满足 \(R^2=\left(\frac a2\right)^2+\left(\frac a2\right)^2+\left(\frac a2\right)^2=\frac{3a^2}{4}\)，球面面积为 \(4\pi R^2=3\pi a^2\).
\end{solution}
\end{problem}

\section{解答题}

\begin{problem}
求函数 \(y=\sin^2 x+2\sin x\cos x+3\cos^2 x\) 的最小值，并写出使函数 \(y\) 取最小值的 \(x\) 的集合.

\begin{solution}
由 \(\sin^2x=\frac{1-\cos2x}{2}\)、\(\cos^2x=\frac{1+\cos2x}{2}\) 及 \(2\sin x\cos x=\sin2x\)，得
\[
y=2+\cos2x+\sin2x=2+\sqrt2\sin\left(2x+\frac{\pi}{4}\right)\geq 2-\sqrt2
\]
当且仅当 \(\sin\left(2x+\frac{\pi}{4}\right)=-1\) 时取等号，即 \(2x+\frac{\pi}{4}=-\frac{\pi}{2}+2k\pi\)，所以 \(x=k\pi-\frac{3\pi}{8}\)，其中 \(k\in\Z\). 因此最小值为 \(2-\sqrt2\)，使函数取最小值的 \(x\) 的集合为 \(\left\{x\mid x=k\pi-\frac{3\pi}{8},\ k\in\Z\right\}\).
\end{solution}
\end{problem}

\begin{problem}
已知复数 \(z=1+\i\)，求复数 \(\frac{z^2-3z+6}{z+1}\) 的模和辐角的主值.

\begin{solution}
因为 \(z=1+\i\)，所以 \(z^2=2\i\)，从而 \(z^2-3z+6=3-\i\)，且 \(z+1=2+\i\). 因此
\[
\frac{z^2-3z+6}{z+1}=\frac{3-\i}{2+\i}=\frac{(3-\i)(2-\i)}{(2+\i)(2-\i)}=1-\i
\]
故该复数的模为 \(\sqrt2\). 因为 \(1-\i\) 对应的点在第四象限，且其辐角的主值取 \([0,2\pi)\) 内的值，所以主值为 \(\frac{7\pi}{4}\).
\end{solution}
\end{problem}

\begin{problem}
已知 \(ABCD\) 是边长为 \(4\) 的正方形，\(E\)，\(F\) 分别是 \(AB\)，\(AD\) 的中点，\(GC\) 垂直于 \(ABCD\) 所在的平面，且 \(GC=2\)，求点 \(B\) 到平面 \(EFG\) 的距离.

\begin{solution}
建立直角坐标系，使正方形所在平面为 \(z=0\)，取 \(A=(0,0,0)\)、\(B=(4,0,0)\)、\(C=(4,4,0)\)、\(D=(0,4,0)\). 由题意可取 \(E=(2,0,0)\)、\(F=(0,2,0)\)、\(G=(4,4,2)\). 于是 \(\overrightarrow{EF}=(-2,2,0)\)，\(\overrightarrow{EG}=(2,4,2)\)，平面 \(EFG\) 的一个法向量为 \((1,1,-3)\)，其方程为 \(x+y-3z-2=0\). 因此点 \(B=(4,0,0)\) 到该平面的距离为
\[
d=\frac{\lvert4+0-3\times0-2\rvert}{\sqrt{1^2+1^2+(-3)^2}}=\frac{2}{\sqrt{11}}=\frac{2\sqrt{11}}{11}
\]
\end{solution}
\end{problem}

\begin{problem}
根据函数单调性的定义，证明函数 \(f(x)=-x^3+1\) 在 \(\left(-\infty,+\infty\right)\) 上是减函数.

\begin{solution}
任取 \(x_1\)，\(x_2\in\left(-\infty,+\infty\right)\)，且 \(x_1<x_2\). 则
\[
f(x_2)-f(x_1)=-(x_2^3-x_1^3)=-(x_2-x_1)(x_1^2+x_1x_2+x_2^2)
\]
其中 \(x_2-x_1>0\). 另一方面，
\[
x_1^2+x_1x_2+x_2^2=\left(x_1+\frac{x_2}{2}\right)^2+\frac34x_2^2\geqslant0
\]
且等号成立当且仅当 \(x_1=x_2=0\)，这与 \(x_1<x_2\) 矛盾，所以 \(x_1^2+x_1x_2+x_2^2>0\). 因此 \(f(x_2)-f(x_1)<0\)，即 \(f(x_1)>f(x_2)\). 根据函数单调性的定义，函数 \(f(x)=-x^3+1\) 在 \(\left(-\infty,+\infty\right)\) 上是减函数.
\end{solution}
\end{problem}

\begin{problem}
已知 \(n\) 为自然数，实数 \(a>1\)，解关于 \(x\) 的不等式：\(\log_a x-4\log_{a^2} x+12\log_{a^3} x+\cdots+n(-2)^{n-1}\log_{a^n} x>\frac{1-(-2)^n}{3}\log_a\left(x^2-a\right)\).

\begin{solution}
由定义域知 \(x>0\) 且 \(x^2-a>0\)，所以必须有 \(x>\sqrt a\). 又由换底公式 \(\log_{a^k}x=\frac1k\log_a x\)，左端等于
\[
\sum_{k=1}^{n}k(-2)^{k-1}\log_{a^k}x=\sum_{k=1}^{n}(-2)^{k-1}\log_a x=\frac{1-(-2)^n}{3}\log_a x
\]
因此原不等式化为 \(\frac{1-(-2)^n}{3}\log_a x>\frac{1-(-2)^n}{3}\log_a\left(x^2-a\right)\).

当 \(n\) 为奇数时，\(\frac{1-(-2)^n}{3}>0\). 因为 \(a>1\)，不等式等价于 \(x>x^2-a\)，即 \(x^2-x-a<0\). 记 \(r=\frac{1+\sqrt{1+4a}}{2}\)，则二次式的两个根为 \(\frac{1-\sqrt{1+4a}}{2}<0\)、\(r>0\)，所以不等式的正根区间为 \(0<x<r\). 与定义域交得 \(\sqrt a<x<r\)；并且由 \(r^2=r+a>a\) 知该区间确实非空.

当 \(n\) 为偶数时，\(\frac{1-(-2)^n}{3}<0\)，不等式变为 \(x<x^2-a\)，即 \(x^2-x-a>0\). 因为负根小于 \(0<\sqrt a\)，结合定义域得 \(x>r\). 综上，当 \(n\) 为奇数时，解集为 \(\left\{x\mid\sqrt a<x<\frac{1+\sqrt{1+4a}}{2}\right\}\)；当 \(n\) 为偶数时，解集为 \(\left\{x\mid x>\frac{1+\sqrt{1+4a}}{2}\right\}\).
\end{solution}
\end{problem}

\begin{problem}
双曲线的中心在坐标原点 \(O\)，焦点在 \(x\) 轴上，过双曲线右焦点且斜率为 \(\sqrt{\frac{3}{5}}\) 的直线交双曲线于 \(P\)，\(Q\) 两点，若 \(OP\perp OQ\)，\(|PQ|=4\)，求双曲线的方程.

\begin{solution}
设双曲线方程为 \(\frac{x^2}{a^2}-\frac{y^2}{b^2}=1\)，其中 \(a>0\)、\(b>0\)，焦距的一半为 \(c\)，则 \(c^2=a^2+b^2\). 令 \(m=\sqrt{\frac35}\)，右焦点为 \((c,0)\)，所给直线为 \(y=m(x-c)\). 设它与双曲线的两个交点的横坐标为 \(x_1\)，\(x_2\)，代入双曲线方程并整理，得
\[
\left(5b^2-3a^2\right)x^2+6a^2cx-3a^2c^2-5a^2b^2=0
\]
若 \(5b^2-3a^2=0\)，直线与双曲线的一条渐近线平行，不能有两个交点，故 \(5b^2-3a^2\ne0\). 由根与系数的关系，记 \(S=x_1+x_2\)、\(T=x_1x_2\)，则
\(S=-\frac{6a^2c}{5b^2-3a^2}\)，\(T=-\frac{a^2(3c^2+5b^2)}{5b^2-3a^2}\).

交点坐标为 \(P=(x_1,m(x_1-c))\)、\(Q=(x_2,m(x_2-c))\). 由 \(OP\perp OQ\)，得
\[
0=x_1x_2+m^2(x_1-c)(x_2-c)=T+\frac35(T-cS+c^2)
\]
即 \(8T-3cS+3c^2=0\). 代入 \(S\)，\(T\) 得
\[
\frac{-8a^2(3c^2+5b^2)+18a^2c^2+3c^2(5b^2-3a^2)}{5b^2-3a^2}=0
\]
再用 \(c^2=a^2+b^2\) 化简为 \(3b^4-8a^2b^2-3a^4=0\)，即 \((3b^2+a^2)(b^2-3a^2)=0\). 因为 \(a\)，\(b>0\)，所以 \(b^2=3a^2\)，从而 \(c=2a\).

此时由上面的根与系数关系得 \(S=-a\)、\(T=-\frac{9a^2}{4}\)，故 \((x_2-x_1)^2=S^2-4T=10a^2\). 又因两点在斜率为 \(m\) 的直线上，\(|PQ|^2=(1+m^2)(x_2-x_1)^2=\frac85\cdot10a^2=16a^2\). 由 \(|PQ|=4\)，得 \(a^2=1\)，于是 \(b^2=3\). 所求双曲线方程为 \(\boxed{x^2-\frac{y^2}{3}=1}\).
\end{solution}
\end{problem}
